%%%PAGE 120 rot=0 legibility=good summary=三球(A,B,C)绳连问题：第1次AC碰撞求v_B、再次共线求v_B与绳拉力F、A最大动能与两绳夹角；页面下方约70%空白
%%%BLOCK id=120-01 ch=07 sec=动量与能量综合·多球系统 kind=problem alt=06,03 cont=none y=0.03-0.31
% 页面右上角露出邻页(化学)内容“…和CO2混合…×22.4=6.72L”，不属于本页，已忽略；页面下方为空白横线(背面透出淡淡反字，已忽略)
\begin{problem}
%FIG p120_a 0.025 0.035 0.148 0.128
\notefig[0.25]{figures/p120_a.png}
% 图中标注：A、B、C 三球在同一水平线上，B 处向上箭头 $v_0$，间距 $L$、$L$

① AC第1次碰\quad 求$v_B$

② 再次三球共线\ $v_B$、绳中拉力$F$

③ 运动过程中A最大动能和两绳夹角
\begin{solution}
① $mv_0=3mv_{\text{共}}$\quad $v_{\text{共}}=\dfrac{v_0}{3}$

② \[
\left\{\begin{aligned}
&mv_0=mv_B+2mv_A\\
&\frac{1}{2}mv_0^2=\frac{1}{2}mv_B^2+\frac{1}{2}m\left(v_A^2+v_A^2\right)
\end{aligned}\right.
\]
\[
\left\{\begin{aligned}
&v_B=-\frac{1}{3}v_0\\
&v_A=\frac{2}{3}v_0
\end{aligned}\right.
\qquad
F=\frac{m v_{\text{相}}^2}{L}=\frac{mv_0^2}{L}
\]
AC相对B\quad $v_{\text{相}}=v_0$

③
%FIG p120_b 0.65 0.118 0.83 0.17
\notefig[0.25]{figures/p120_b.png}
% 图中标注：$v_A$（左上箭头）、$v_C$（右上箭头）、45°，两处直角记号，顶点处有夹角记号

沿绳$v_B$为0，故AC均无分量\quad $E_{kA\max}$\quad $v_B=0$

$v_A$\ $v_C$垂绳\quad $\theta=90^\circ$
\[ 2m_Av\sin 45^\circ=mv_0 \]% CHECK: 下标 A 的位置疑应在 v 上(2mv_A sin45°=mv_0)
\[ v_A=v_C=\frac{\sqrt{2}}{2}v_0 \]
\[ E_{kA}=E_{kC}=\frac{mv_0^2}{4} \]
$\theta=90^\circ$ % 此 θ=90° 在页面中部偏右、字较大，疑为 ③ 的最终结论
\end{solution}
\end{problem}
%%%END
%%%PAGEEND 120
